% Paquets généraux
\documentclass[a4paper,12pt,titlepage,twoside]{article}
\usepackage[utf8]{inputenc}
\usepackage[T1]{fontenc}
\usepackage[french]{babel}
\usepackage[gen]{eurosym}
%\usepackage[dvips]{graphicx}
\usepackage{fancyhdr}
\usepackage{pdfpages}
\usepackage{multido}
\usepackage{hyperref}
\usepackage{numprint}
\nplpadding{2}
\usepackage{pstricks}
%\usepackage{aeguill}

\newcommand{\id}{11}
\newcommand{\nom}{Analyse de la structure d'un système asservi}
\newcommand{\sequence}{02}
\newcommand{\num}{02}
\newcommand{\type}{C}
\newcommand{\descrip}{Définition et structure des systèmes asservis, insertion dans les chaînes d'information et d'énergie, les critères de qualité d'un système asservi}
\newcommand{\competences}{A3-C11: Structure des systèmes asservis \\ &  Mod2-C4: Systèmes linéaires continus invariants asservis \\ &  Mod2-C20: Modélisation des systèmes asservis}
\newcommand{\nbcomp}{3}
\newcommand{\systemes}{}
\newcommand{\ilot}{0}

\newcommand\twodigits[1]{%
   \ifnum#1<10 0#1\else #1\fi
}

\newcommand{\auteurun}{Renaud Costadoat}
\newcommand{\auteurdeux}{Françoise Puig}
\newcommand{\auteurs}{Costadoat - Puig}
\newcommand{\institute}{Lycée Dorian}
\newcommand{\scriptpython}{\href{https://raw.githubusercontent.com/Costadoat/S\sequence/master/\type\num\space\nom/Ilot_\twodigits{\ilot} \systemes/\sequence-\type\num-I\twodigits{\ilot}.py}{script python }}

\usepackage{color}
\usepackage{xcolor}
\usepackage{colortbl}
\usepackage{helvet}
\renewcommand{\familydefault}{\sfdefault}
\usepackage{amsfonts}
\usepackage{amsmath}
%\usepackage{xspace}
\usepackage{varioref}
\usepackage{tabularx}
%\usepackage{floatflt}
\usepackage{graphics}
\usepackage{wrapfig}
\usepackage{textcomp}
\usepackage{tikz}
\usepackage{wrapfig}
\usepackage{gensymb}
\usepackage[european]{circuitikz}
\usetikzlibrary{babel}
\usepackage{ifthen}
\usepackage{cancel}
\usepackage{etoolbox}
\usepackage{multirow}
%\usepackage{boxedminipage}
\definecolor{gris25}{gray}{0.75}
\definecolor{analyser}{RGB}{196,189,151}
\definecolor{analyserf}{RGB}{73,69,41}
\definecolor{modeliser}{RGB}{252,213,180}
\definecolor{modeliserf}{RGB}{226,107,10}
\definecolor{resoudre}{RGB}{216,228,188}
\definecolor{resoudref}{RGB}{118,147,60}
\definecolor{experimenter}{RGB}{230,184,183}
\definecolor{experimenterf}{RGB}{150,54,52}
\definecolor{concevoir}{RGB}{141,180,226}
\definecolor{concevoirf}{RGB}{15,36,62}
\definecolor{realiser}{RGB}{184,204,228}
\definecolor{realiserf}{RGB}{54,96,146}
\definecolor{communiquer}{RGB}{204,192,218}
\definecolor{communiquerf}{RGB}{96,73,122	}
\definecolor{bleu}{RGB}{18,33,98}
\definecolor{bleutf}{RGB}{8,53,154}
\definecolor{bleuf}{RGB}{42,94,171}
\definecolor{bleuc}{RGB}{231,239,247}
\definecolor{rougef}{RGB}{185,18,27}
\definecolor{rougec}{RGB}{255,188,204}%255,230,231
\definecolor{vertf}{RGB}{103,126,82}
\definecolor{vertc}{RGB}{220,255,191}
\definecolor{forestgreen}{rgb}{0.13,0.54,0.13}
\definecolor{blcr}{rgb}{0.59,0.69,0.84}
\definecolor{blfr}{rgb}{0.32,0.51,0.75}
\definecolor{orfr}{rgb}{0.90,0.42,0.15}
\definecolor{orcr}{rgb}{0.90,0.65,0.50}
\definecolor{orangef}{rgb}{0.89,0.56,0.16}
\definecolor{orange}{rgb}{0.58,0.35,0.063}
\definecolor{orangec}{rgb}{0.94,0.76,0.56}
\definecolor{rcorrect}{rgb}{0.6,0,0}
\definecolor{sequence}{rgb}{0.75,0.75,0.75}
\definecolor{competences}{rgb}{0.61,0.73,0.35}
\definecolor{grisf}{HTML}{222222}
\definecolor{grisc}{HTML}{636363}
\definecolor{normal}{HTML}{4087c4}
\definecolor{info}{HTML}{5bc0de}
\definecolor{success}{RGB}{92,184,92}
\definecolor{warning}{RGB}{240,173,78}
\definecolor{danger}{RGB}{217,83,79}
\hypersetup{                    % parametrage des hyperliens
    colorlinks=true,                % colorise les liens
    breaklinks=true,                % permet les retours à la ligne pour les liens trop longs
    urlcolor= blfr,                 % couleur des hyperliens
    linkcolor= orange,                % couleur des liens internes aux documents (index, figures, tableaux, equations,...)
    citecolor= forestgreen                % couleur des liens vers les references bibliographiques
    }

% Mise en page
\pagestyle{fancy}

\setlength{\hoffset}{-18pt}

\setlength{\oddsidemargin}{0pt} 	% Marge gauche sur pages impaire2s
\setlength{\evensidemargin}{0pt} 	% Marge gauche sur pages paires
\setlength{\marginparwidth}{00pt} 	% Largeur de note dans la marge
\setlength{\headwidth}{481pt} 	 	% Largeur de la zone de tête (17cm)
\setlength{\textwidth}{481pt} 	 	% Largeu\textbf{r de la zone de texte (17cm)
\setlength{\voffset}{-18pt} 		% Bon pour DOS
\setlength{\marginparsep}{7pt}	 	% Séparation de la marge
\setlength{\topmargin}{-30pt} 		% Pas de marge en haut
\setlength{\headheight}{45pt} 		% Haut de page
\setlength{\headsep}{20pt} 			% Entre le haut de page et le texte
\setlength{\footskip}{30pt} 		% Bas de\textbf{ page + séparation
\setlength{\textheight}{700pt} 		% Hauteur de l'icone zone de texte (25cm)
\setlength\fboxrule{1 pt}
\renewcommand{\baselinestretch}{1}
\setcounter{tocdepth}{1}
\newcommand{\cadre}[2]
{\fbox{
  \begin{minipage}{#1\linewidth}
   \begin{center}
    #2\\
   \end{center}
  \end{minipage}
 }
}

\newcounter{num_quest} \setcounter{num_quest}{0}
\newcounter{num_rep} \setcounter{num_rep}{0}
\newcounter{num_cor} \setcounter{num_cor}{0}

\newcommand{\question}[1]{\refstepcounter{num_quest}\par
~\ \\ \parbox[t][][t]{0.15\linewidth}{\textbf{Question \arabic{num_quest}}}\parbox[t][][t]{0.93\linewidth}{#1}\par
}


\newcommand{\reponse}[1]
{\refstepcounter{num_rep}
\noindent
\rule{\linewidth}{.5pt}
\textbf{Question \arabic{num_rep}:}
\multido{\i=1+1}{#1}{~\ \\}
}

\newcommand{\cor}
{\refstepcounter{num_cor}
\noindent
\rule{\linewidth}{.5pt}
\textbf{Question \arabic{num_cor}:} \\
}

\newcommand{\titre}[1]
{\begin{center}
\cadre{0.8}{\huge #1}
\end{center}
}

%Titres des sections

\newcommand{\titleana}[2]{%
\tikzstyle{titlebox}=[rectangle,inner sep=10pt,inner ysep=10pt,draw,fill=analyser]%
\tikzstyle{title}=[fill=analyser]%
%
\bigskip\noindent\begin{tikzpicture}
\node[titlebox] (box){%
\begin{minipage}{0.94\textwidth}
\textbf{\textcolor{analyserf}{#1}}
\end{minipage}
};
%\draw[color=blue] (box.north west)--(box.north east);
\node[title] at (box.north) {\textcolor{analyserf}{ANALYSER}};
\end{tikzpicture}%
\\}

\newcommand{\titlemod}[2]{%
\tikzstyle{titlebox}=[rectangle,inner sep=10pt,inner ysep=10pt,draw,fill=modeliser]%
\tikzstyle{title}=[fill=modeliser]%
%
\bigskip\noindent\begin{tikzpicture}
\node[titlebox] (box){%
\begin{minipage}{0.94\textwidth}
\textbf{\textcolor{modeliserf}{#1}}
\end{minipage}
};
%\draw[color=blue] (box.north west)--(box.north east);
\node[title] at (box.north) {\textcolor{modeliserf}{MODELISER}};
\end{tikzpicture}%
\\}

\newcommand{\titleres}[2]{%
\tikzstyle{titlebox}=[rectangle,inner sep=10pt,inner ysep=10pt,draw,fill=resoudre]%
\tikzstyle{title}=[fill=resoudre]%
%
\bigskip\noindent\begin{tikzpicture}
\node[titlebox] (box){%
\begin{minipage}{0.94\textwidth}
\textbf{\textcolor{resoudref}{#1}}
\end{minipage}
};
%\draw[color=blue] (box.north west)--(box.north east);
\node[title] at (box.north) {\textcolor{resoudref}{RESOUDRE}};
\end{tikzpicture}%
\\}

\newcommand{\titleexp}[2]{%
\tikzstyle{titlebox}=[rectangle,inner sep=10pt,inner ysep=10pt,draw,fill=experimenter]%
\tikzstyle{title}=[fill=experimenter]%
%
\bigskip\noindent\begin{tikzpicture}
\node[titlebox] (box){%
\begin{minipage}{0.94\textwidth}
\textbf{\textcolor{experimenterf}{#1}}
\end{minipage}
};
%\draw[color=blue] (box.north west)--(box.north east);
\node[title] at (box.north) {\textcolor{experimenterf}{EXPERIMENTER}};
\end{tikzpicture}%
\\}

\newcommand{\titlecon}[2]{%
\tikzstyle{titlebox}=[rectangle,inner sep=10pt,inner ysep=10pt,draw,fill=concevoir]%
\tikzstyle{title}=[fill=concevoir]%
%
\bigskip\noindent\begin{tikzpicture}
\node[titlebox] (box){%
\begin{minipage}{0.94\textwidth}
\textbf{\textcolor{concevoirf}{#1}}
\end{minipage}
};
%\draw[color=blue] (box.north west)--(box.north east);
\node[title] at (box.north) {\textcolor{concevoirf}{CONCEVOIR}};
\end{tikzpicture}%
\\}

\newcommand{\titlerea}[2]{%
\tikzstyle{titlebox}=[rectangle,inner sep=10pt,inner ysep=10pt,draw,fill=realiser]%
\tikzstyle{title}=[fill=realiser]%
%
\bigskip\noindent\begin{tikzpicture}
\node[titlebox] (box){%
\begin{minipage}{0.94\textwidth}
\textbf{\textcolor{realiserf}{#1}}
\end{minipage}
};
%\draw[color=blue] (box.north west)--(box.north east);
\node[title] at (box.north) {\textcolor{realiserf}{REALISER}};
\end{tikzpicture}%
\\}

\newcommand{\titlecom}[2]{%
\tikzstyle{titlebox}=[rectangle,inner sep=10pt,inner ysep=10pt,draw,fill=communiquer]%
\tikzstyle{title}=[fill=communiquer]%
%
\bigskip\noindent\begin{tikzpicture}
\node[titlebox] (box){%
\begin{minipage}{0.94\textwidth}
\textbf{\textcolor{communiquerf}{#1}}
\end{minipage}
};
%\draw[color=blue] (box.north west)--(box.north east);
\node[title] at (box.north) {\textcolor{communiquerf}{COMMUNIQUER}};
\end{tikzpicture}%
\\}


\newcommand{\prob}[2]{%
\tikzstyle{titlebox}=[rounded corners=10,inner sep=10pt,inner ysep=10pt,draw,fill=bleuf]%
%\tikzstyle{title}=[]%
%
\bigskip\noindent\begin{tikzpicture}
\node[titlebox] (box){%
\begin{minipage}{0.94\textwidth}
\begin{minipage}{0.15\linewidth}
\includegraphics[width=0.8\linewidth]{../../../img/questionmark}
\end{minipage}
\begin{minipage}{0.8\linewidth}
\textbf{\textcolor{white}{
 Problématique du TP:\\ \noindent\rule{13cm}{1pt} \\ \\#1}}
\end{minipage}
\end{minipage}
};
%\draw[color=blue] (box.north west)--(box.north east);
%\node[title] at (box.north west) {\textcolor{white}{\textbf{Problématique}}};
\end{tikzpicture}\bigskip\\~\ \\%
}


% En tête et pied de page
\lhead{\nom}
\rhead{\includegraphics[width=2cm]{../../../img/logo}}
\lfoot{\auteurs}
\cfoot{Page \thepage}

\fancypagestyle{correction}{%
  \fancyhf{}
  \lhead{\colorbox{danger}{\begin{minipage}{0.65\paperwidth} \textcolor{white}{\textbf{Correction}} \end{minipage}} }
  \rhead{\includegraphics[width=2cm]{../../../img/logo}}
  \lfoot{Renaud Costadoat}
  \rfoot{\colorbox{danger}{\begin{minipage}{0.6\paperwidth} \begin{flushright}\textcolor{white}{\textbf{Correction}}\end{flushright} \end{minipage}} }}

\renewcommand{\footrulewidth}{0.4pt}

\usepackage{eso-pic}
\newcommand{\BackgroundPic}{%
\put(0,0){%
\parbox[b][\paperheight]{\paperwidth}{%
\vfill
\begin{center}
\hspace{0.5cm}\vspace{0.5cm}
\includegraphics[width=\paperwidth,height=\paperheight,%
keepaspectratio]{../../../img/fond3}%
\end{center}
\vfill
}}}

\newcommand{\BackgroundPicdeux}{%
\put(25,-30){%
\parbox[b][\paperheight]{\paperwidth}{%
\vfill
\begin{center}
\includegraphics[width=\paperwidth,height=\paperheight,%
keepaspectratio]{../../../img/fond4}%
\end{center}
\vfill
}}}

\newcommand{\activite}[1]{\newpage \cleardoublepage
\renewcommand\arraystretch{2.5}\setcounter{section}{0}\setcounter{subsection}{0}
\ifthenelse{\equal{#1}{1}}{\begin{center}\begin{tabular}{|c|cc p{3cm}|}
\hline
\rowcolor{violt} \supportu & Act & Contenu & Compétences \\\hline
 \multirow{3}{*}{\includegraphics[width=0.2\linewidth]{\supportimageu}} & \begin{LARGE}\textbf{1}\end{LARGE} & \activiteun & \competenceun \\
 & \multicolumn{3}{l|}{Nom: .........................................................................}\\
 & \multicolumn{3}{l|}{Prénom: ....................................................................}\\\hline
\end{tabular}\end{center}}{}
\ifthenelse{\equal{#1}{2}}{\begin{center}\begin{tabular}{|c|cc p{3cm}|}
\hline
\end{tabular}\end{center}}{}}

\begin{document}

\pagestyle{empty}

\vspace*{-3\baselineskip}

\AddToShipoutPicture*{\BackgroundPic}

\ifdef{\auteurdeux}{\begin{tabular}{>{\columncolor{gray!00}}m{.33\linewidth} m{.27\linewidth} >{\columncolor{gray!00}}m{.3\linewidth}}
Séquence \sequence \ - \type \num \ - Îlot \numprint{\ilot} &  \multirow{3}{*}{\hspace{1cm}\includegraphics[height=1.5cm]{../../../img/logo}} &  \begin{flushright} \multirow{4}{*}{\hspace{1cm}\includegraphics[height=4cm]{img/qrcode}}\end{flushright}\\
 \textbf{\institute} \\
 \auteurun\\
 \auteurdeux
\end{tabular}}{\begin{tabular}{>{\columncolor{gray!00}}m{.3\linewidth} m{.3\linewidth} >{\columncolor{gray!00}}m{.3\linewidth}}
Séquence \sequence \ - \type \num \ - Îlot \numprint{\ilot}  &  \multirow{3}{*}{\hspace{1cm}\includegraphics[height=1.5cm]{../../../img/logo}} &  \begin{flushright} \multirow{4}{*}{\hspace{1cm}\includegraphics[height=4cm]{img/qrcode}}\end{flushright}\\
 \textbf{\institute} \\
 \auteurun
\end{tabular}}

\vspace{1cm}

\begin{center}\colorbox{white}{\huge{\nom}}\end{center}

\vspace{2cm}


\begin{center}\colorbox{white!20}{\includegraphics[height=5cm]{/home/renaud/Documents/Renaud/GitHub/django_education/django_education/static/systemes/\imageun}}\end{center}

\vspace{2cm}



\begin{tabular}{p{.15\linewidth} >{\columncolor{white}}p{.8\linewidth}}
    \rowcolor{gray!20}
    Référence & S\sequence \ - \type\num \ - I\numprint{\ilot} \\
    Compétences & \competences \\
 	\rowcolor{gray!20}
    Description & \descrip \\
    Système & \systemes
  \end{tabular}

\newpage

\AddToShipoutPicture{\BackgroundPicdeux}

\pagestyle{fancy}


\ifdef{\public}{\cleardoublepage}{}

\prob{Modéliser un Système Linéaire Continu et Invariant à l'aide d'un schéma bloc} 

\titleexp{Identifier les blocs sur le système}

L'objectif de ce TP est de modéliser le berceau automatique Mobycrea par le schéma blocs
suivant : (L'étude portera sur la régulation de vitesse sur l'axe horizontal du Mobycrea).

\begin{center}
 \includegraphics[width=0.95\linewidth]{img/mobycrea}
\end{center}

L'objectif de cette première partie va être de déterminer les fonctions de transfert manquantes
dans la description du système. $Ka$ est un gain pur qui permet d'adapter le signal d'entrée. $Kc$ est
un gain qui symbolise un correcteur. On prendra et $Kc=1$.

\begin{center}
 \includegraphics[width=0.95\linewidth]{img/chaine_energie}
\end{center}

La chaine d'énergie du mouvement horizontal est représentée ci-dessus, le moteur est alimenté
par un commutateur (transistor), qui reçoit les consignes de la chaîne d'information. Le moteur
entraine le siège en translation par l'intermédiaire du réducteur de vitesse et du système bielle
manivelle.

\paragraph{Question 1:} Fonction de transfert du transistor : nous allons considérer qu'il se comporte comme un gain pur $Kh$. Écrire alors sa fonction de transfert $H_1(p)$.


\paragraph{Question 2:} Fonction de transfert du moteur à courant continu : vous avez déjà déterminé la fonction de transfert du moteur à courant continu (sans la charge). Écrire la fonction de transfert $H_m(p)$.

La figure suivante permet de comprendre comment on obtient la rotation de la grande roue (en rose) qui transmettra le mouvement horizontal par le biais d'une bielle (rouge).

\begin{center}
 \includegraphics[width=0.95\linewidth]{img/engrenages}
\end{center}

La carte de commande est informée de la vitesse du moteur grâce à un capteur optique associé à un disque à secteurs, l'ensemble constitue un codeur incrémental. Le codeur incrémental génère des impulsions en tout ou rien (pulses) contrairement à un codeur absolu qui aurait donné la position angulaire.

\paragraph{Question 3:} Fonction de transfert du système de réduction de vitesse : Déterminez le nombre d'impulsions reçues pour une rotation de la bielle (la grande roue comporte 90 dents, la petite 24 et la vis un filet, le codeur à 6 secteurs). 

Par quel moyen peut-on obtenir la vitesse angulaire de la roue à partir des impulsions générées par le capteur ? Déterminer
alors $Ke$, $Kcp$ et $Ka$.

\ifdef{\public}{\cleardoublepage}{}

\titleana{Structure de l'asservissement du système}

L'objectif de cette partie est de déterminer le comportement du système en analysant la structure de l'asservissement.
~\

\paragraph{Question 7:} Déterminer la FTBO du système $FTBO(p)$.

\paragraph{Question 8:} Déterminer la FTBF du système $FTBF(p)$.

\paragraph{Question 9:} Déterminer le gain de la FTBF, le temps de réponse à 5\%, et les constantes du système ($\tau$, $\xi$ et $\omega_0$) et l'écart statique.

\paragraph{Question 10:} Vérifier les données trouvées à la question précédente grâce au tracé du système réel.
 
\ifdef{\public}{\cleardoublepage}{}

\titlemod{Simulation du comportement du modèle}

Le logiciel \textbf{Scilab} permet de tracer la réponse temporelle d'une fonction de transfert donnée.

Pour cela, il suffit de lancer le logiciel et d'aller dans le module \textbf{Xcos}.

Dans le dossier \textbf{CPGE} du navigateur de palettes, vous trouverez, par exemple:
\begin{itemize}
 \item une \textit{entrée}: STEP\_FUNCTION,
 \item un \textit{Opérateur linéaire}: CLR, vous modifierez sa fonction de transfert afin d'obtenir ce que vous souhaitez observer,
 \item une \textit{sortie}: SCOPE,
 \item un \textit{outil d'analyse}: REP\_TEMP.
\end{itemize}

Faire glisser ces blocs sur une page vierge du module xcos et cliquer sur la flèche permettant de lancer la simulation.

\paragraph{Question 11:} Effectuer le tracé du schéma bloc du système sur Scilab.

\paragraph{Question 12:} Tracer les réponses temporelles obtenues durant les activités précédentes afin de vérifier le modèle choisi.

\ifdef{\public}{\end{document}}{}

\newpage

\pagestyle{correction}\setcounter{section}{0}

Modélisation
\begin{center}
$H(p)=\dfrac{F_c(p)}{U_m(p)}=\dfrac{\dfrac{K_m}{R_m.R_p.r}}{1+\dfrac{K_e.K_m}{R_m.K_c.R_p^2.r^2}.p+\dfrac{R_m.J}{R_m.K_c.R_p^2.r^2}.p^2}$
\end{center}

\end{document}